\scanpage{126}
\mspage{-III- \handwritten{7} -}
\begin{flushright}n° \handwritten{383}\end{flushright}
quotient un groupe essentiel dans l'espace euclidien $E/W'$.

On dit que $G$ est \uline{réductible} si $E$ est somme directe de deux sous-espaces orthogonaux $E',E''$ distincts de $E$ et stables pour ${}^vG$. Alors tout élément $s$ de $G$ définit, par passage au quotient, des automorphismes affines $s'$ de $E/E''$ (identifié à $E'$) et $s''$ de $E/E'$ (identifié à $E''$) ; l'ensemble des $s'$ (resp. $s''$) est un sous-groupe $G'$ (resp. $G''$) du groupe affine de $E'$ (resp. $E''$). Si $H$ est un hyperplan de $E$ tel que $s_H\in G$, on voit aisément que, ou bien $E'$ est contenu dans la direction ${}^vH$ de $H$, ou bien $E'$ contient la droite orthogonale à $H$ ; dans le premier cas l'automorphisme $s'$ de $E'$ défini par $s_H$ est l'identité, et dans le second $s'$ est la réflexion par rapport à l'hyperplan $H\handwritten{\cap}E'$ de $E'$. Il s'ensuit que $G'$ et $G''$ sont engendrés par des réflexions. D'autre part, l'application $s\handwritten{\rightsquigarrow}(s',s'')$ de $G$ dans $G'\handwritten{\times}G''$ est évidemment un isomorphisme.

On dit que $G$ est \uline{irréductible} s'il n'est pas réductible. Notant $\underline{H}$ l'ensemble des hyperplans tels que $s_H\in G$, ceci revient à dire qu'il n'existe pas de partition de $\underline{H}$ en deux sous-ensembles tels que tout hyperplan du premier soit perpendiculaire à tout hyperplan du second. Dans le cas général il résulte de ce qu'on vient de voir que $E$ s'identifie khanoniquement à un produit d'espaces euclidiens $E^{(1)},\ldots,E^{(q)}$, et $G$ à un produit de groupes $G^{(1)},\ldots,G^{(q)}$ opérant dans ces espaces, de telle sorte que chaque $G^{(i)}$ soit irréductible.

Nous nous fixerons désormais un ensemble $\underline{H}$ d'hyperplans (affines) de $E$ tel que

1) les réflexions $s_H$ $(H\in\underline{H})$ engendrent $G$,

2) les hyperplans de $\underline{H}$ sont permutés par $G$.

Comme au n°1, notons que, pour $g\in G$, on a $s_{g(H)}=gs_Hg^{-1}$. Nous verrons plus loin que l'ensemble des $s_H$ $(H\in\underline{H})$ est l'ensemble de toutes les réflexions de $G$.

PROPOSITION 4.- \uline{La famille d'hyperplans} $\underline{H}$ \uline{est localement finie}.

Raisonnons par l'absurde. Soit $x\in E$ tel que tout voisinage $U$ de $x$ rencontre une infinité d'hyperplans de $\underline{H}$. Alors $g(U)\handwritten{\cap}U\handwritten{\ne\varnothing}$ pour une infinité de réflexions affines
% archive sheet number: 125
